There were $n$ coders from Finland and $n$ coders from Sweden in a programming contest. It turned out that after the contest, each coder had a distinct score.
Your task is to find the $k$-th highest score in the contest.
To do this, you can ask questions: you can choose a country (Finland or Sweden) and an integer $i$ and you will be told the $i$-th highest score for the chosen country.

\vspace{0.3cm}\noindent\textbf{Interaction}

This is an interactive problem. Your code will interact with the grader using standard input and output. You should start by reading two integers $n$ and $k$.
On your turn, you can print one of the following:

\textbullet\ "$\mathrm{F}\ i$", where $1 \le i \le n$: ask the $i$-th highest score for Finland.
\textbullet\ "$\mathrm{S}\ i$", where $1 \le i \le n$: ask the $i$-th highest score for Sweden.
\textbullet\ "$!\ s$": report that the $k$-th highest score is $s$. Your program must terminate after this.

Each line should be followed by a line break. You must make sure the output gets flushed after printing each line.

\vspace{0.3cm}\noindent\textbf{Constraints}

\textbullet\ $1 \le n \le 10^5$
\textbullet\ $1 \le k \le 2n$
\textbullet\ each score is between $1$ and $10^9$
\textbullet\ you can ask at most $100$ queries of the first two types in total